\section{The original invariant covariance and its conductor-pole filtration}\label{the-original-invariant-covariance-and-its-conductor-pole-filtration}

21 September 2026. Independent continuation of CG29--32. All coefficient rows below are the actual invariant rows surviving the complete low-jet constraint. No replacement of their metric, their source projection or their original period is made.

This calculation proves two quantitative statements. First, the sum of the positive logarithmic eigenvalues of the original invariant covariance, relative to its specified original coefficient form, is \(O_A(q\log(q+2))=o(kq)\). Second, the remaining graph-covariance changes in CG32 can be restricted to the \textbf{actual invariant rows that vanish at the specified conductor poles}, with error \(O_A(q\sqrt{k}\log^2(q+2))=o(kq)\). The removed rows are controlled by proved metric bounds before their number is used. On the retained rows the entire conductor covariance has a quadratic-form upper bound \(\exp(o(q))\) in that coefficient form. Its small eigenvalues and its comparison with the outer-ideal covariance remain to be evaluated.

\subsection{1. Every row and every analytic function comes from the original data}\label{every-row-and-every-analytic-function-comes-from-the-original-data}

Retain KF1--55, CG1--32, the complete OCF invariant construction, and \[
N\in\{q-1,q,2q-1,2q\},\quad L=N-q,
D=N-v,\quad r=8k-32-v+s_k.
\tag{CGR1}
\] Let \(Z\in\mathbb C^{r\times q}\) contain the actual coefficient rows obtained in CG29--30. It is supported on the OCF strip \(E_k^{\rm strip}\), has full row rank, and each row annihilates the entire degree-\(<v\) polynomial space. It is a linear combination of the complete columns \(J_R=S_kZ_k\), followed by the exact elimination of \(\operatorname{im}B_{\mathcal L}\). Its reference coefficient form is \[
R_Z=ZZ^*>0.
\tag{CGR2}
\] This is only a fixed reference used to compare matrices. Every actual Gamma metric below remains unchanged.

For a row \(z=(z_\alpha)\) in this original row space define \[
F_z(w)=\sum_{\alpha=1}^qz_\alpha e^{\beta_\alpha w},\qquad
G_z(w)=F_z(w)/E_A(w).
\tag{CGR3}
\] Since the row kills all degree-\(<v\) polynomials, \(F_z^{(j)}(0)=0\) for \(j<v\). Thus \(G_z\) is holomorphic at zero. The numerator is the original invariant biform restricted to the original exponential curve. OCF17 supplies an invariant lift also for the strip representative: subtracting its conductor part stays inside the invariant image. No arbitrary exponential numerator has been inserted.

For every polynomial \(p\), finite differentiation gives the exact duality \[
[p(\partial_w)F_z(w)]_{w=0}
=[(\mathcal T_Ap)(\partial_w)G_z(w)]_{w=0}.
\tag{CGR4}
\] For a single monomial this is Leibniz's formula for \(F_z=E_AG_z\); linearity proves the general case. Consequently the complete target observation is \[
\ell_z(\eta)=
[\eta(\partial_w)\chi'(\partial_w)G_z(w)]_{w=0}.
\tag{CGR5}
\] Indeed KF8--10 give \(\mathcal T_A\mathcal R(\chi'\eta)=\chi'\eta\), and CGR4 gives precisely CG29's coefficient row on every low polynomial \(\eta\). The extension CGR5 annihilates the \textbf{entire} relation: \[
\ell_z(\mathcal U_Ah)
=[(\mathcal T_A(\chi h))(\partial_w)G_z]_{0}
=[(\chi h)(\partial_w)F_z]_{0}=0,
\tag{CGR6}
\] because every exponent \(\beta_\alpha\) is a root of the full original \(\chi\). Thus it is the actual quotient observation at every cutoff, not an extension chosen for a norm estimate.

\subsection{2. The exact original Gamma covariance, with its projection retained}\label{the-exact-original-gamma-covariance-with-its-projection-retained}

Let \(A=A_{D,1}\) be the complete WCF6 matrix from the upper-center Gamma frame \(\phi_{v,c},\ldots,\phi_{N,c}\) to the lower-center frame \(\phi_{0,c'},\ldots,\phi_{D,c'}\). Its entries and every phase are unchanged. Let \(B_F\in\mathbb C^{r\times(D+1)}\) have entries \[
(B_F)_{j,n}=[\phi_{v+n,c}(\partial_w)F_{z_j}(w)]_{0}.
\tag{CGR7}
\] In the lower-center orthonormal coefficient space let \(P_I\) be the orthogonal projection onto the exact polynomial subspace \[
I_N=\chi'\mathcal P'_{L+g}\subset\mathcal P'_D.
\tag{CGR8}
\] The target observation covariance is exactly \[
\boxed{C_{R,N}^U=B_FA^{-1}P_I(A^{-1})^*B_F^*.}
\tag{CGR9}
\] To prove it, CGR4 on the source Gamma basis gives the lower functional row \(B_FA^{-1}\). Its restriction to \(I_N\) is its composition with the orthogonal projection \(P_I\); the squared Gram of those restricted rows is CGR9. Multiplication by \(\chi'\) is exactly the isometry from the weighted target source into \(I_N\). CGR6 shows that taking the complete quotient afterward leaves that observation covariance unchanged. The projection \(P_I\), of codimension \(q'\), is present throughout.

Set \(R_y=k\sqrt{\delta^2+\gamma^2}\) and retain the exact finite root-evaluation bound from KF33, \[
V_N=\frac{qe^2}{M_\sigma}(N+1)^{3/2}(2N+1)^{R_y+1}.
\tag{CGR10}
\] Then \(B_FB_F^*\preceq V_NR_Z\). This follows by writing the rows as the original coefficient rows times the complete root-evaluation matrix; discarding the first \(v\) source columns does not increase its norm. The bound itself follows from the generating function on radius \(N/(N+1)\), as proved in KF33 and CG23.

With \(M_D,J_D\) exactly as in CG23, the WCF19d exterior estimate gives, for each \(1\le j\le r\), \[
\prod_{i=1}^j\lambda_i
\le V_N^j M_D^{2(D+1-j)}e^{2J_D},
\tag{CGR11}
\] where \(\lambda_1\ge\cdots\ge\lambda_r>0\) are the eigenvalues of \(R_Z^{-1/2}C_{R,N}^UR_Z^{-1/2}\). In fact factor its rectangular Gram matrix as \(R_Z^{-1/2}B_FA^{-1}P_I\), apply the exterior norm to each factor, and use that an orthogonal projection contracts every exterior norm. The product of the largest \(j\) squared singular values is the displayed left side. Taking the largest partial logarithmic sum proves \[
\boxed{\sum_{i=1}^r\log^+\lambda_i
\le r\log^+V_N+2(D+1)\log M_D+2J_D
=O_A(q\log(q+2))=o(kq).}
\tag{CGR12}
\] Here \(r=O(k)\), \(\log V_N=O_A(k\log(q+2))\), and the WCF quantities have their already proved orders. This controls every positively amplified covariance direction jointly. It neither bounds the negative logarithms by the same expression nor assumes the retained projection is an identity.

\subsection{3. Finite global bounds needed before eliminating pole directions}\label{finite-global-bounds-needed-before-eliminating-pole-directions}

These bounds will pay for the actual eliminated directions, including their angles. Let \(a_*^{\rm piv}\ne0\) be the original leading conductor coefficient selected in OCF5, and put \[
C_A=1+A_1/|a_*^{\rm piv}|,\qquad
B_{\rm val}=qM_\sigma
\left(\frac{2q+R_y}{2\delta}\right)^{2(q-1)}.
\tag{CGR13}
\] Each OCF6 coefficient elimination has coefficient \(\ell^1\) operator norm at most \(C_A\): it subtracts a scalar coefficient times the full shifted conductor polynomial divided by its actual pivot. There are \(q'\) steps. Hence its exact strip reduction satisfies \(\|N_k\|_2\le\sqrt q C_A^{q'}\). For any original strip row \(z\), and any conductor coefficient vector \(M_Au\), \[
\|z-M_Au\|_2\ge\|z\|_2/\|N_k\|_2.
\tag{CGR14}
\] This follows from \(N_kM_A=0\) and \(N_kz=z\) in the identical strip positions. It is a quantitative consequence of the actual full conductor elimination, not an assumption about the angle of the invariant image.

The complete source covariance in root-value coordinates is at least \(B_{\rm val}^{-1}I_q\), by the explicit Lagrange interpolation estimate KF34. The lower bound can also be read directly from the inverse norm of the interpolation polynomials: each root separation has modulus at least \(2\delta\), and successive multiplication by \(y-\omega\) has Gamma norm at most \(2q+R_y\) on the relevant degrees. The sum of their squared norms is exactly bounded by CGR13.

The orthogonal complement of \(I_N\) is the span of evaluations at every original lower root. Apply CGR4 to these evaluation rows. WCF21--22 gives their images under \(A\) as exactly the complete upper conductor rows. All their first \(v\) Gamma coefficients, like those of \(F_z\), vanish. Consequently the distance defining the left side of CGR9 is at least the distance of the original source row to the conductor image divided by \(\|A\|\). Using CGR14 and the full source lower bound proves \[
l_{U,N}R_Z\preceq C_{R,N}^U\preceq u_{U,N}R_Z,
\quad
l_{U,N}=\frac1{M_D^2 B_{\rm val}q C_A^{2q'}},\qquad
u_{U,N}=V_N M_D^{2D}e^{2J_D}.
\tag{CGR15}
\] For the upper bound use CGR11 at \(j=1\). In particular \(\log^+u_{U,N}+\log^+(1/l_{U,N})=O_A(q\log(q+2))\). No lower eigenvalue is asserted to be of order one.

For clarity, the corresponding outer-ideal covariance has finite bounds in the \textbf{same} reference form, not a newly chosen observation frame. At \(N_0=q-1\), CG24 proves \(C_{R,N_0}^d=C_{R,N_0}^U\). The outer observation is fixed polynomial remainder followed by \(A_R\), so enlarging its source adds nonnegative rank-one covariance terms. Thus \(C_{R,N}^d\succeq l_{U,N_0}R_Z\) at all later cutoffs.

Here are explicit upper constants. Write \(M=L+g\) for the target source degree, \(c'=k/2-4\), and keep CG6's finite domain. Put \[
\begin{split}
R_d&=\max\{1,|c'|+q/r_*\},\qquad t_d=2R_d,\\
B_{\rm rem}(M)&=g2^g(1+R_d)^g\sqrt{M+1}\,t_d^M,\\
h_-(M)&=\frac{C_Le^{-\pi(q+1)}(q^2-R_y^2)^{q'}}
{(M+1)^2(2M+1)64^M(1+|c'|+q)^{2M}},\\
h_+^0&=gM_\sigma(4q+2+R_y+|c'|)^{2q}.
\end{split}\tag{CGR16}
\] Ordinary remainder by \(d_A/a_A\) on the coefficient vectors of \(1,S',\ldots,S'^M\) has norm at most \(B_{\rm rem}(M)\). To prove it, use the exact Cauchy remainder formula on \(|S'|=t_d\). Its denominator has modulus at least \((t_d/2)^g\); its monic coefficient sum is at most \((1+R_d)^g\). The coefficient sum of its divided difference is at most \(g t_d^{g-1}(1+R_d)^g\). Contour length supplies \(t_d\), and a source polynomial has boundary maximum at most \(\sqrt{M+1}t_d^M\) times its coefficient norm. This proves the bound, including repeated roots, since the contour remainder retains all jets.

The complete source Gram \(H_M\) in this coefficient frame obeys \(H_M\succeq h_-(M)I\). On \(q\le y\le q+1\), paired lower roots give \(|\chi'(c'+iy)|^2\ge(q^2-R_y^2)^{q'}\), and CG4's Gamma density bound contributes \(C_Le^{-\pi(q+1)}\). Expanding on \([0,1]\) in shifted Legendre polynomials, as in CG17, bounds the coefficient norm after the substitution \(S'=c'+i(q+x)\) by \((M+1)\sqrt{2M+1}\,8^M(1+|c'|+q)^M\) times the interval \(L^2\) norm. Squaring gives \(h_-(M)\). The low Gram \(H_{g-1}\) has norm at most \(h_+^0\): multiply its monomial columns by all \(q'\) original root factors and use CG11's upper bound on each, then sum their squared norms. All these products have degree at most \(q\).

At the first cutoff \(C_{R,N_0}^U=A_RH_{g-1}^{-1}A_R^*\), hence \(A_RA_R^*\preceq h_+^0u_{U,N_0}R_Z\). Applying the complete remainder bound to the high source covariance now proves \[
l_dR_Z\preceq C_{R,N}^d\preceq u_{d,N}R_Z,
\quad l_d=l_{U,N_0},\qquad
u_{d,N}=\frac{B_{\rm rem}(M)^2h_+^0u_{U,N_0}}{h_-(M)}.
\tag{CGR17}
\] Every logarithmic bound in CGR17 has size \(O_A(q\log(q+2))\), since \(M\le q+g=O(q)\), \(g=O(k)\), and \(R_d=O_A(q)\). These deliberately coarse global constants are used only on the few rows removed next.

\subsection{4. An exact growing family of conductor-pole constraints}\label{an-exact-growing-family-of-conductor-pole-constraints}

Retain the full WCF function and the original center difference: \[
w(t)=-i\arctan t,\quad
g_A(t)=t^{-v}e^{-4w(t)}E_A(w(t)),\quad
g_A(0)=g_0=(-i)^v\mu_v/v!\ne0.
\tag{CGR18}
\] The powers and logarithm are their analytic branches at \(t=0\), on the unit disk. For the original row \(z\), its upper generating function and the lower generating function of \(G_z\) satisfy \[
\begin{split}
f_z(t)&=t^{-v}(1+t^2)^{-1/4}e^{-cw(t)}F_z(w(t)),\\
H_z(t)&=(1+t^2)^{-1/4}e^{-c'w(t)}G_z(w(t)),\qquad
f_z(t)=g_A(t)H_z(t).
\end{split}\tag{CGR19}
\] The exact coefficients of \(H_z\) are \([p_n((\partial_w-c')/i)G_z(w)]_0/n!\). This follows by applying the original Gamma generating function to the linear functional. The factors \(c,c'\), \(i\), and \(e^{-4w}\) all remain. The apparent singularity of \(f_z\) at zero is removable because the first \(v\) source coefficients vanish.

Put \[
\eta=k^{-1/2},\quad \rho=1-2\eta,\quad
R_1=1-\eta,\quad R_2=1-\eta/2,
\tag{CGR20}
\] which belong to \((0,1)\) for the original \(k\ge9\). Define \[
M_g=\max\left\{1,|g_0|,
2^vA_1e^{2\pi\gamma}
\left(\frac{1+R_2}{1-R_2}\right)^{4\delta}\right\},
\quad J_g=\log(M_g/|g_0|)\ge0.
\tag{CGR21}
\] WCF7--8 bounds \(|g_A|\) by \(M_g\) on \(|t|\le R_2\). Let its zeros in \(|t|<R_1\) be \(\zeta\), with their actual multiplicities \(m_\zeta\), and let \(H_k=\sum m_\zeta\). Jensen's finite disk formula gives \[
H_k\le\frac{J_g}{\log(R_2/R_1)}\le\frac{2J_g}{\eta}
=O_A(\sqrt{k}\log(k+2)).
\tag{CGR22}
\] Indeed every such zero contributes at least \(\log(R_2/R_1)\) to the Jensen sum at \(R_2\), while the boundary mean is at most \(\log M_g\). The formula follows by factoring its zeros and applying the mean-value formula to the logarithm of the remaining nonzero analytic factor; this is the same full proof used in OR9. The second inequality follows from \(\log(R_2/R_1)\ge(R_2-R_1)/R_2\ge\eta/2\).

Define a subspace \textbf{inside the original invariant row space} by \[
\mathcal W_k^0=
\{z:\ F_z^{(j)}(w(\zeta))=0
\text{ for every }\zeta\text{ above and }0\le j<m_\zeta\}.
\tag{CGR23}
\] Its exact codimension \(h_k\) is the rank of the displayed finite matrix of actual values \(\beta_\alpha^j e^{\beta_\alpha w(\zeta)}\) after multiplication by \(Z\); in particular \(h_k\le H_k\). Every derivative and point is specified by the original conductor and invariant rows. Since \(w'(t)=-i/(1+t^2)\ne0\), these constraints say exactly that the numerator \(f_z\) cancels the indicated zeros of \(g_A\) with their full multiplicities. There is no unspecified generic rank.

\subsection{5. The complete projected covariance on those rows has sub-q upper cost}\label{the-complete-projected-covariance-on-those-rows-has-sub-q-upper-cost}

Let \(B(t)\) be the finite disk Blaschke product of the zeros in CGR22, with factors \(R_1(t-\zeta)/(R_1^2-\bar\zeta t)\), repeated with their multiplicities. Its boundary modulus on \(|t|=R_1\) is one and \(|B(0)|\le1\). Both \(f_z/B\) and \(g_A/B\) are holomorphic there for \(z\in\mathcal W_k^0\), and the second has no interior zero. The maximum principle gives \(\max|f_z/B|\le\max_{|t|=R_1}|f_z(t)|\). Also \(u(t)=\log(M_g/|g_A(t)/B(t)|)\) is nonnegative harmonic in that disk, with \(u(0)\le J_g\). Its Poisson integral gives \[
u(t)\le\frac{R_1+\rho}{R_1-\rho}u(0)
\le\frac{2J_g}{\eta}\qquad(|t|\le\rho).
\tag{CGR24}
\] For the inequality, the Poisson kernel on any intermediate circle of radius \(R< R_1\) is at most \((R+|t|)/(R-|t|)\); its average of \(u\) is \(u(0)\). Let \(R\) tend to \(R_1\). Boundary zeros of \(g_A\), if any, do not affect this interior argument or the maximum principle.

The numerator bound retaining its original coefficient norm is \[
\max_{|t|=R_1}|f_z(t)|
\le F_k\|z\|_2,\qquad
F_k=R_1^{-v}(1-R_1^2)^{-1/4}\sqrt q
\exp\{R_y\operatorname{arctanh}R_1\}.
\tag{CGR25}
\] In fact the original center factor cancels \(e^{cw}\) in each original exponential; the remaining exponent has modulus at most \(R_y|w|\), and \(|w(t)|\le\operatorname{arctanh}R_1\). Cauchy--Schwarz on all \(q\) unchanged coefficients proves CGR25. Using \(M_g\ge1\), CGR24--25 therefore bound the full lower generating function by \(\max_{|t|\le\rho}|H_z(t)|\le F_ke^{2J_g/\eta}\|z\|_2\). Its coefficient of degree \(n\) has modulus at most this quantity times \(\rho^{-n}\). The normalized Gamma functional coefficient multiplies that coefficient by \(\rho_n/\sqrt{M_\sigma}\), where \(\rho_n^2=n!/(1/2)_n\le2n+1\). Summing every source coefficient through \(D\), and then retaining the original orthogonal projection, proves \[
\boxed{C_{R,N}^U\big|_{\mathcal W_k^0}
\preceq e^{\varepsilon_{k,N}}R_Z\big|_{\mathcal W_k^0},}
\quad
\varepsilon_{k,N}=
\log\frac{(D+1)(2D+1)}{M_\sigma}+2\log F_k
+\frac{4J_g}{\eta}+2D\log(1/\rho).
\tag{CGR26}
\] An upper bound using \(\max(0,\varepsilon_{k,N})\) is available if a nonnegative allowance is desired. Explicitly \[
\varepsilon_{k,N}
=O_A(k^{3/2}+k\log(k+2)+\sqrt{k}\log(k+2))=o(q).
\tag{CGR27}
\] Here \(D\le2q\), \(\log(1/\rho)=O(k^{-1/2})\), \(\operatorname{arctanh}R_1=O(\log k)\), and \(J_g=O_A(\log k)\). The proof never replaces the projected norm by an equality with the full norm; the full norm is used only for its valid upper bound.

\subsection{6. The actual allocation loses only o(kq) when these poles are removed}\label{the-actual-allocation-loses-only-okq-when-these-poles-are-removed}

Choose a basis orthonormal for the fixed coefficient form \(R_Z\), whose first \(r-h_k\) rows span the exact subspace CGR23. Use that same basis at both high cutoffs and for both \(U,d\). In its block decomposition the determinant equals the determinant of the first restriction times the determinant of its attained Schur complement. Denote the covariance change on the first restriction by \(\Delta_{R,N}^0\).

Congruence and attained minimization transport CGR15 and CGR17 to those Schur complements: every eigenvalue for \(U\) lies in \([l_{U,N},u_{U,N}]\), and every one for \(d\) lies in \([l_d,u_{d,N}]\). Their rank is exactly \(h_k\). Therefore \[
\boxed{|\Delta_{R,N}-\Delta_{R,N}^0|
\le h_k\max\left\{
\log^+\frac{u_{U,N}}{l_d},
\log^+\frac{u_{d,N}}{l_{U,N}}\right\}
=O_A(q\sqrt{k}\log^2(q+2))=o(kq).}
\tag{CGR28}
\] This is where the number of pole directions is used. Their complete metric costs were bounded first, so no unproved angle estimate or deletion by rank is concealed in CGR28. The coefficient basis is only an exact fixed change of frame and its determinant cancels in each covariance ratio.

Combining with the already evaluated full-volume error CG32 gives \[
\boxed{\mathcal G_K
=-\Delta_{R,2q-1}^0-\Delta_{R,2q}^0+o_A(kq).}
\tag{CGR29}
\] For a finite bound, add CG32's displayed finite right side to the two finite right sides of CGR28. Thus the remaining leading allocation is now localized to an explicitly defined, original invariant row space with every selected conductor pole canceled. It still has dimension \(r-O_A(\sqrt{k}\log(k+2))\). CGR26--27 control its upper covariance quadratic form. The lower spectrum and the signed comparison with \(C_{R,N}^d\) on this space are not evaluated here; an upper bound alone cannot assign their determinant ratio or an RH conclusion.

\subsection{Sources and exact reading coverage}\label{sources-and-exact-reading-coverage}

OCF1--26, \emph{An exact conductor frame and a fixed invariant-generator recursion}, was read completely for this continuation. CGR3 and CGR14 use its exact strip representatives, conductor columns and invariant certificates. The \href{https://github.com/KokunoYumeto/zeta-function-research-reader/blob/5992ad50c0f2a06921f1fcaded7b4e38097c0739/workbenches/splitzero-tandem/continuations/20260920-original-kernel-mixed-determinants/providers/ORIGINAL_CONDUCTOR_STRIP_FRAME.tex\#L27}{complete pinned proof file} is byte-verified in the root public-source receipt.

KF1--55 and CG1--32 were read completely. This continuation retains their actual \(A_R\), full low constraint, frame and source. Their complete source must accompany publication before a live link is claimed.

WCF4--13, WCF19a--f and WCF20--22 are used from the \href{https://github.com/KokunoYumeto/zeta-function-research-reader/blob/52274804051e325c8a6ef69c17b3958cc2caba0a/workbenches/splitzero-tandem/continuations/20260920-gamma-and-finite-derivative/providers/WEIGHTED_CONDUCTOR_FORWARD.tex}{pinned complete proof file}, whose byte identity was checked in the CG source. The cited portions were read directly. The Gamma generating function is NIST DLMF \href{https://dlmf.nist.gov/18.23.E7}{18.23.7}, with recurrence \href{https://dlmf.nist.gov/18.22.E8}{18.22.8}, by T. H. Koornwinder, R. Wong, R. Koekoek, R. F. Swarttouw and W. P. Reinhardt, version 1.2.8. Its original TeX encodings are retained in \nolinkurl{graph_primary_sources/}. The generating function is applied explicitly in CGR19 with both original centers and their exact phase. Jensen's zero count is the classical formula of J. L. W. V. Jensen, \emph{Sur un nouvel et important théorème de la théorie des fonctions} (1899), \href{https://doi.org/10.1007/BF02417878}{DOI 10.1007/BF02417878}; the finite proof used here is included in OR9 and recalled in CGR22. The disk-product and harmonic inequalities needed in CGR24 are proved above rather than supplied as missing analytic assumptions.

The independent checker \nolinkurl{check_invariant_covariance.py} passed 96 exact tests of CGR3--9, with v=0 and v=1 finite-shift fixtures, both retained Gamma centers, the physical complex phase, the complete relation minimum and the original source projection. Deleting that projection or changing the phase fails the respective negative control. These finite checks do not replace the analytic proofs of CGR10--29.

The exact earlier proof files are \href{https://github.com/KokunoYumeto/zeta-function-research-reader/blob/0fd693c844aad9117e30ec447c59864a28af10f3/workbenches/splitzero-tandem/continuations/20260921-native-spectral-real-pair/003/ACTUAL_KERNEL_FILTRATION_AND_GRAPH.tex\#L12}{KF1--55}, \href{https://github.com/KokunoYumeto/zeta-function-research-reader/blob/0fd693c844aad9117e30ec447c59864a28af10f3/workbenches/splitzero-tandem/continuations/20260921-native-spectral-real-pair/003/COMPLETE_CONDUCTOR_GRAPH_METRIC.tex}{CG1--32}, and \href{https://github.com/KokunoYumeto/zeta-function-research-reader/blob/0fd693c844aad9117e30ec447c59864a28af10f3/workbenches/splitzero-tandem/continuations/20260921-native-spectral-real-pair/003/COMPLETE_CONDUCTOR_ROOT_GEOMETRY.tex\#L38}{OR1--25}. CGR and CGS are reproduced completely in this edition.
